% UzBLiMP paper draft (2026-10-28, benchmark v1.1). Use with the ACL template on Overleaf:
% https://www.overleaf.com/latex/templates/association-for-computational-linguistics-acl-conference/jvxskxpnznfj
% See paper/README.md. Compiler: pdfLaTeX (the template default).
% All tables in tables/ are generated by scripts/paper_tables.py; numbers in the text come from them
% or from results/eval_v1.1.md.
\documentclass[11pt]{article}

% [review]: anonymous, line-numbered submission version.
% For the arXiv preprint use \usepackage[preprint]{acl}; for camera-ready use \usepackage{acl}.
\usepackage[review]{acl}

\usepackage{times}
\usepackage{latexsym}
\usepackage[T2A,T1]{fontenc}  % T2A for the Cyrillic examples
\usepackage[utf8]{inputenc}
\usepackage{microtype}
\usepackage{inconsolata}
\usepackage{graphicx}
\usepackage{booktabs}
\usepackage{xcolor}

% Uzbek Latin apostrophes: oʻ/gʻ (U+02BB) and the tutuq belgisi (U+02BC)
\DeclareUnicodeCharacter{02BB}{\textquoteleft}
\DeclareUnicodeCharacter{02BC}{\textquoteright}
\newcommand{\cyr}[1]{{\fontencoding{T2A}\selectfont #1}}
\newcommand{\todo}[1]{\textcolor{red}{[TODO: #1]}}
\newcommand{\ungr}[1]{*\textit{#1}}

\title{UzBLiMP: Uzbek Minimal Pairs in Latin, Cyrillic and the Reformed 2026 Alphabet}

% Authors are hidden in [review] mode. TODO: agree authorship with Javokhir Matniyazov
% and Russell Scheinberg before adding them (notes/decisions.md, decisions 11 and 13).
\author{Kattabek Norbaev \\
  Independent researcher, Uzbekistan \\
  \texttt{channel0101k@gmail.com}}

\begin{document}
\maketitle

\begin{abstract}
We introduce UzBLiMP, a benchmark of 2,400 linguistic minimal pairs for Uzbek covering 16 grammatical phenomena, 150 pairs each. Every pair is derived from an attested corpus sentence by one controlled edit and was reviewed by native speakers. Because Uzbek is written in two scripts and its Latin alphabet was reformed in 2026, we release the benchmark in current Latin, Cyrillic and the reformed Latin alphabet. Uzbek-only language models score 95.7--98.6\% in the script they were trained on, while a small multilingual model is at chance because it prefers shorter sentences. Two Uzbek models trained separately on Latin and Cyrillic text score nearly the same (95.7\% vs.\ 94.6\%), but a model reading a script it was not trained on loses 7--22 points: its input fragments into twice as many tokens and it falls back on the same preference for shorter sentences. On the reformed spelling, a Latin-trained model loses 11 points; most of the loss comes from pairs in which only the grammatical sentence contains a new letter, where the model drops below chance by preferring the familiar spelling. Minimal pairs that differ in orthographic familiarity or length therefore measure more than grammar. \todo{update with the larger models and the human baseline}
\end{abstract}

\section{Introduction}

Benchmarks of linguistic minimal pairs test whether a language model assigns higher probability to a grammatical sentence than to a minimally different ungrammatical one \citep{warstadt-etal-2020-blimp}. Because the two sentences differ in a single controlled place, the comparison isolates one piece of grammatical knowledge without asking the model to follow instructions, which makes the method usable for small models and for languages without instruction data. Following BLiMP for English, such benchmarks now exist for several languages, including the Turkic languages Turkish \citep{basar2025turblimp} and Tatar \citep{saetov2026tatblimp}. Uzbek, a Turkic language with more than 30 million speakers \todo{cite a source for the speaker count}, has so far been covered only by 50 subject--verb agreement pairs in the massively multilingual MultiBLiMP \citep{jumelet2025multiblimp}.

Uzbek also offers a natural experiment in writing systems. In the last century it has been written in Arabic script, in Latin from 1928, in Cyrillic from 1940, and again in Latin since the 1990s \citep{fierman1991}; Latin and Cyrillic text coexist today. In 2026 the Latin alphabet was reformed again, replacing oʻ, gʻ, sh and ch with ö, ğ, ş and ç \todo{cite the official decision}. Text in three orthographies therefore exists side by side, and models trained before the reform will meet the new spelling. This raises a question with practical consequences: \emph{does what a language model knows about grammar survive a change of alphabet?}

We make four contributions:
\begin{itemize}
\item \textbf{UzBLiMP}, 2,400 minimal pairs for 16 phenomena (12 following the TurBLiMP inventory, 4 specific to Uzbek), built from attested corpus sentences with one controlled edit and reviewed by native speakers (\S\ref{sec:benchmark}).
\item \textbf{Three aligned alphabet versions}: current Latin, Cyrillic, and, to our knowledge, the first evaluation resource in the reformed 2026 Latin alphabet (\S\ref{sec:alphabets}).
\item \textbf{An analysis of script dependence}: grammar knowledge is learnable from either script but does not transfer between them, and the failure has an identifiable mechanism, token fragmentation followed by a fallback on sentence length (\S\ref{sec:results}).
\item \textbf{Two methodological cautions} for minimal-pair evaluation: differences in orthographic familiarity and in length between the two members can masquerade as grammatical knowledge or its absence (\S\ref{sec:discussion}).
\end{itemize}

\section{Related work}

\paragraph{Minimal-pair benchmarks.} BLiMP \citep{warstadt-etal-2020-blimp} introduced 67 paradigms of template-generated English pairs and established accuracy on forced choice by sentence probability as the standard metric. TurBLiMP \citep{basar2025turblimp} adapted the design to Turkish with 16 phenomena, combining hand-written and hand-corrected pairs with automatically generated ones, and found that length and subword-count differences between the members explain little of the models' performance. TatBLiMP \citep{saetov2026tatblimp} combined real Tatar sentences with sentences composed from corpus-verified vocabulary, required ungrammatical forms to be real word forms (the plausibility principle), and used byte-normalised accuracy as its main metric. MultiBLiMP \citep{jumelet2025multiblimp} covers agreement phenomena in many languages, including 50 Uzbek pairs. UzBLiMP follows the phenomenon inventory of TurBLiMP and the plausibility principle of TatBLiMP, but builds every pair from an attested corpus sentence. \todo{add SLING, RuBLiMP, CLiMP and other language-specific benchmarks after checking the references}

\paragraph{Uzbek language models and resources.} Goldfish \citep{chang2024goldfish} provides small monolingual models for hundreds of languages, including separate Uzbek models trained on comparable amounts of Latin-script and Cyrillic-script text, which makes them a natural testbed for script effects. tahrirchi-bert \citep{tahrirchi-bert} is an Uzbek masked language model. UzLiB \citep{tahrirchi-uzlib} evaluates Uzbek linguistic knowledge with multiple-choice questions; UzBLiMP complements it with probability-based minimal pairs that do not require instruction following. Our pipeline uses TurkicNLP \citep{turkicnlp} for transliteration and part-of-speech tags.

\paragraph{Grammar knowledge and prompting.} \citet{scheinberg-etal-2025-explain} show that stating the relevant grammatical rule in the prompt improves acceptability judgments in English, Chinese and Russian, especially for smaller models. Whether such explanations help when the learned pattern is weak, as in a low-resource language or an unfamiliar script, is an open question that UzBLiMP makes testable. \todo{if the grammar-prompting experiment is included, expand here}

\section{Uzbek and its alphabets}
\label{sec:alphabets-bg}

Uzbek is an agglutinative Turkic language with SOV order, case and possessive suffixes on nouns, person agreement on verbs, and pervasive suffixation for tense, aspect, mood and nominalisation \citep{sjoberg1963,bodrogligeti2003}. Several of its grammatical regularities are morphological: the dative suffix has allomorphs conditioned by the stem-final consonant, stem-final \textit{q}/\textit{k} voice before vowel-initial suffixes, and a possessed noun agrees in person with its possessor. These properties make Uzbek minimal pairs often differ by a single suffix.

The current Latin alphabet (1995) writes the rounded back vowel as oʻ and the voiced uvular fricative as gʻ, and uses the digraphs sh and ch. The Cyrillic alphabet, still widely used, has dedicated letters for these sounds (e.g.\ \cyr{ў} for oʻ, \cyr{ш} for sh, \cyr{ч} for ch). The 2026 reform keeps the Latin script but replaces oʻ, gʻ, sh and ch by single letters (ö, ğ, ş, ç). Converting between these systems is deterministic for most of the vocabulary but not for Russian loanwords, whose Cyrillic spelling (e.g.\ the soft sign in \cyr{фильм}) is not recoverable from Latin (\S\ref{sec:alphabets}).

\section{Building UzBLiMP}
\label{sec:benchmark}

\subsection{Phenomena}
Table~\ref{tab:phenomena} lists the 16 phenomena. Twelve follow the TurBLiMP inventory where the contrast exists in Uzbek: anaphor agreement and binding with the reflexive \textit{oʻz-}, lexical case of transitive verbs (\textit{qoʻrqmoq} `to fear' takes the ablative), the dative of ditransitive recipients, suspended affixation in coordination, nominalisation, person agreement, passives with an agent phrase, relative clauses, the quantifier \textit{har bir} `every', and NPI licensing of \textit{hech}-words by negation. Some Turkish designs were changed where Uzbek differs; for example, an agent phrase with \textit{tomonidan} requires a passive verb. Phenomena whose Uzbek counterpart was unclear in our native-speaker checks (ellipsis, scrambling, impersonal passives, determiners, islands) were dropped. TurBLiMP's irregular-forms slot has no Uzbek counterpart and is filled by a morphophonological alternation: the voicing of stem-final \textit{q}/\textit{k} before vowel-initial possessive suffixes (\textit{oʻrtogʻim}, \ungr{oʻrtoqim}). Four phenomena are specific to Uzbek: no plural after numerals, case government of postpositions (\textit{keyin} and \textit{beri} take the ablative), possessive agreement, and the dative allomorphs \textit{-qa}/\textit{-ka} after stem-final \textit{q}/\textit{k} (\textit{qishloqqa}, \ungr{qishloqga}). Appendix~\ref{app:phenomena} describes each phenomenon and its edit.

\begin{table*}[t]
\centering\small
\begin{tabular}{@{}lp{7.7cm}p{5.5cm}@{}}
\toprule
Phenomenon & Minimal pair & Translation of the grammatical sentence \\
\midrule
\multicolumn{3}{@{}l}{\textit{Following the TurBLiMP inventory}} \\
Anaphor agreement & Biz \textbf{oʻzimizni}/\ungr{oʻzimni} noqulay his qilardik. & `We used to feel uncomfortable.' \\
Lexical case & Faqat men \textbf{undan}/\ungr{uni} qoʻrqaman. & `Only I am afraid of him.' \\
Ditransitive dative & Men \textbf{unga}/\ungr{uni} yaralarimni koʻrsatdim. & `I showed him my wounds.' \\
Binding & Biroq men hozirgacha \textbf{oʻzimni}/\ungr{meni} koyiyman. & `But to this day I blame myself.' \\
Stem alternation (q/k) & Bu shimpanzening \textbf{yuragi}/\ungr{yuraki} edi. & `This was the chimpanzee's heart.' \\
NPI \textit{hech} & Gapimga hech kim \textbf{ishonmadi}/\ungr{ishondi}. & `Nobody believed what I said.' \\
Suspended affixation & Pul \textbf{topdim}/\ungr{topgan} va xotinimga yubordim. & `I earned money and sent it to my wife.' \\
Nominalization & Bu tajribalarni qanday \textbf{tushuntirish}/\ungr{tushuntirgan} mumkin? & `How can these experiments be explained?' \\
Person agreement & Avvaliga men bir oz \textbf{hayajonlandim}/\ungr{hayajonlandik}. & `At first I was a little nervous.' \\
Passives & Asar tanqidchilar tomonidan yuqori \textbf{baholangan}/\ungr{baholagan}. & `The work was rated highly by critics.' \\
Relative clauses & Polda \textbf{singan}/\ungr{sindi} shisha butilka yotibdi. & `A broken glass bottle is lying on the floor.' \\
Quantifiers (\textit{har bir}) & Har bir \textbf{mahalla}/\ungr{mahallalar} tashkiliy tuzilmaga ega. & `Every neighbourhood has an organisational structure.' \\
\midrule
\multicolumn{3}{@{}l}{\textit{Specific to Uzbek}} \\
Numeral + singular & Oilada toʻrtta \textbf{farzand}/\ungr{farzandlar} bor. & `There are four children in the family.' \\
Postposition case & Buni \textbf{bolaligimdan}/\ungr{bolaligimga} beri bilaman. & `I have known this since my childhood.' \\
Possessive agreement & Ogohlantirib qoʻyish mening \textbf{vazifam}/\ungr{vazifang}. & `It is my duty to give a warning.' \\
Dative allomorphy & Qayiq ohistagina \textbf{qirgʻoqqa}/\ungr{qirgʻoqga} keldi. & `The boat came gently to the shore.' \\
\bottomrule
\end{tabular}
\caption{The 16 phenomena of UzBLiMP with one pair each (150 pairs per phenomenon). The grammatical form of the edited word is in bold and the ungrammatical form follows the slash. Interlinear glosses are in Appendix~\ref{app:phenomena}.}
\label{tab:phenomena}
\end{table*}


\subsection{Generating candidate pairs}
Each pair consists of an attested sentence and a copy with exactly one word changed by a phenomenon-specific operation, such as replacing a dative pronoun by its accusative form or a passive verb by its active form. Following TatBLiMP, the changed word must itself be a plausible Uzbek word form: it must occur at least twice in a frequency lexicon of 988,950 word forms, except for the two morphophonology phenomena, whose ungrammatical forms are by definition non-words. Sentences come from two sources:
\begin{itemize}
\item an Uzbek Latin corpus of 64 million words (7.2 million sentences, 131,000 documents of books, Wikipedia and legal texts), processed by rule-based generators for 13 phenomena;
\item 4,500 rule-generated pairs for 8 phenomena built by a team member from the \textit{uz-books} collection \citep{tahrirchi-uzbooks}, of which 1,388 passed our filters.
\end{itemize}
All sentences are single sentences of 3--16 words without digits, quotation marks, dialogue or proper names after the first word. Generators use TurkicNLP part-of-speech tags to check the edited word (e.g.\ that the word after a numeral is a noun) and skip known false positives, such as function words that only look possessive (\textit{ham} `also' parsed as \textit{ha}+1\textsc{sg}) or accusative pronouns that could be read as the object of a nearby participle. Every pair in the benchmark is derived from an attested corpus sentence. Appendix~\ref{app:generation} gives the filters and guards.

\subsection{Native-speaker review}
\label{sec:review}
Of 3,800 candidates, 3,729 were reviewed and 3,267 (87.6\%) accepted; acceptance per phenomenon ranges from 72\% to 95\% (Table~\ref{tab:phenomena}). Each candidate was judged first by the first author, a native speaker of Uzbek, and the judgments were then checked with three other native speakers. Of the rejections, 71\% concern the source sentence (OCR errors, typos, fragments) rather than the edit; 13\% are pairs whose ``ungrammatical'' member has a grammatical reading, concentrated in relative clauses, where a finite verb before a noun can be read as the end of a clause; 12\% involve dialectal or colloquial forms (Table~\ref{tab:reasons} in the appendix). For each phenomenon, 150 accepted pairs were sampled with a fixed seed; the remaining 867 accepted pairs form a reserve. \todo{expert review by linguists at a local university (planned)} \todo{inter-annotator agreement \citep{cohen1960} on a blind sample judged independently by another native speaker}

\subsection{Benchmark statistics}
Table~\ref{tab:stats} summarises the benchmark. Grammatical sentences have 8.8 words on average. 1,438 pairs come from our corpus and 962 from uz-books; five phenomena draw on both sources. Phenomena built on closed-class items have little lexical variety at the edited position (4 distinct forms for binding, 5 for anaphor agreement, 6 for ditransitive recipients), whereas open-class phenomena vary widely (146 forms for possessive agreement). The two members of a pair usually differ in length by a few bytes; the sign of this difference matters for models that prefer shorter sentences (\S\ref{sec:length}): the ungrammatical member is shorter on average for binding ($-4.9$ bytes, \textit{oʻzimni} vs.\ \textit{meni}) and NPI licensing ($-2.0$, the negation suffix is removed), and longer for numerals and \textit{har bir} ($+3.0$, the plural suffix is added).

\begin{table}[t]
\centering\small
\setlength{\tabcolsep}{3.6pt}
\begin{tabular}{@{}lrrrrr@{}}
\toprule
Phenomenon & Books & Words & Forms & $\Delta$B & Acc. \\
\midrule
Anaphor agreement & 150 & 6.3 & 5 & $-$0.8 & 95 \\
Lexical case & 0 & 7.8 & 98 & $-$0.4 & 86 \\
Ditransitive dative & 78 & 8.2 & 6 & $-$0.7 & 88 \\
Binding & 150 & 7.2 & 4 & $-$4.9 & 88 \\
Stem alternation & 0 & 8.6 & 57 & $-$0.4 & 90 \\
NPI \textit{hech} & 0 & 7.6 & 108 & $-$2.0 & 92 \\
Suspended affixation & 0 & 10.7 & 95 & 0.0 & 93 \\
Nominalization & 122 & 9.6 & 110 & +0.2 & 72 \\
Subject agreement & 150 & 6.9 & 127 & $-$0.1 & 91 \\
Passives & 81 & 10.6 & 105 & $-$1.5 & 84 \\
Relative clauses & 107 & 10.7 & 114 & $-$1.0 & 74 \\
Quantifiers & 122 & 11.6 & 130 & +3.0 & 77 \\
Numeral + singular & 0 & 8.9 & 136 & +3.0 & 93 \\
Postposition case & 0 & 8.7 & 90 & $-$1.0 & 91 \\
Possessive agreement & 0 & 7.9 & 146 & +0.5 & 92 \\
Dative allomorphy & 0 & 8.8 & 93 & 0.0 & 92 \\
\midrule
All & 960 & 8.8 & & $-$0.4 & 88 \\
\bottomrule
\end{tabular}
\caption{Benchmark statistics per phenomenon (150 pairs each). \emph{Books}: pairs whose sentence comes from UzBooks, the others come from our corpus; \emph{Words}: mean length of the grammatical sentence; \emph{Forms}: distinct grammatical word forms at the edited position; $\Delta$\emph{B}: mean length of the ungrammatical minus the grammatical sentence in UTF-8 bytes; \emph{Acc.}: share (\%) of reviewed candidates accepted in the native-speaker review.}
\label{tab:stats}
\end{table}


\subsection{Three alphabets}
\label{sec:alphabets}
\paragraph{Cyrillic.} We convert every pair with TurkicNLP's transliterator and a set of fixes for errors we found in it: \textit{t}+\textit{s} across morpheme boundaries is \cyr{тс}, not \cyr{ц} (\textit{ketsa} \textrightarrow{} \cyr{кетса}); a tutuq between \textit{s} and \textit{h} only separates the letters; \textit{ye} after a consonant is a Russian loan written with a hard or soft sign; \textit{n}+\textit{gʻ} is two letters, not the digraph \textit{ng} plus a hard sign; Roman numerals are kept; and month names and a small list of loan stems receive the soft sign (\textit{film} \textrightarrow{} \cyr{фильм}). After conversion no Latin letters remain and every pair still differs in one word. A native speaker checked a sample of 10 pairs per phenomenon \todo{80/80 correct for 8 phenomena; the other 80 pending}. Appendix~\ref{app:conversion} lists the rules.

\paragraph{Reformed Latin.} We apply the four replacements of the 2026 reform (oʻ\textrightarrow ö, gʻ\textrightarrow ğ, sh\textrightarrow ş, ch\textrightarrow ç). For each pair we record whether the changed word contains a new letter (1,070 pairs), only other words do (1,208), or nothing changes (122). Among the first group we further distinguish pairs where the new letter appears in both members' changed words (819) from pairs where it appears only in the grammatical member (251; e.g.\ \textit{özimni} vs.\ \ungr{meni}).

\section{Experimental setup}

\paragraph{Models.} We evaluate four models: the Goldfish Uzbek models trained on Latin (uzn\_latn\_100mb) and on Cyrillic text (uzn\_cyrl\_100mb) \citep{chang2024goldfish}; tahrirchi-bert-base, an Uzbek masked LM \citep{tahrirchi-bert}; and the multilingual Qwen2.5-0.5B \citep{qwen2025qwen25}. The two Goldfish models share architecture and training-data size and differ in script, which makes them our main comparison for script effects. Goldfish-latn is not scored on reformed Latin because 82\% of the converted sentences contain a token unknown to its tokenizer; Goldfish-latn reading Cyrillic and Goldfish-cyrl reading Latin are scored despite unknown tokens (28\% and 81\% of sentences) to show what happens when a model meets a script it was not trained on. \todo{add instruction-tuned 7--9B models and a frontier model}

\paragraph{Scoring.} Causal models are scored with the summed log-probability of the sentence using minicons \citep{misra2022minicons}, with a beginning-of-sequence token so that the first word is scored; the masked model is scored with pseudo-log-likelihood \citep{salazar-etal-2020-masked} using the within-word left-to-right variant \citep{kauf-ivanova-2023-better}. A pair is correct when the grammatical sentence scores strictly higher. Following TatBLiMP we also report length-normalised accuracy, comparing log-probability per UTF-8 byte (Appendix~\ref{app:results}). All scores are computed in 32-bit precision, which we found necessary because lower precision made scores depend on batch size. We compare conditions on the same pairs with McNemar's test \citep{mcnemar1947} and report bootstrap 95\% confidence intervals (2,000 resamples).

\section{Results}
\label{sec:results}

\begin{table*}[t]
\centering\small
\begin{tabular}{lrrrrrrrrrr}
\toprule
 & \multicolumn{4}{c}{Latin text} & \multicolumn{4}{c}{Cyrillic text} & \multicolumn{2}{c}{Reformed text} \\
\cmidrule(lr){2-5} \cmidrule(lr){6-9} \cmidrule(lr){10-11}
Phenomenon & GF-L & GF-C & BERT & Qwen & GF-L & GF-C & BERT & Qwen & BERT & Qwen \\
\midrule
Anaphor agreement & 98.7 & 40.7 & 99.3 & 50.7 & 88.0 & 99.3 & 72.0 & 38.0 & 86.7 & 42.0 \\
Lexical case & 98.7 & 45.3 & 100.0 & 31.3 & 98.7 & 98.0 & 84.0 & 28.7 & 87.3 & 28.0 \\
Ditransitive dative & 96.7 & 45.3 & 99.3 & 31.3 & 94.7 & 94.0 & 89.3 & 24.0 & 98.7 & 39.3 \\
Binding & 90.7 & 0.0 & 100.0 & 7.3 & 67.3 & 81.3 & 39.3 & 1.3 & 29.3 & 4.0 \\
Stem alternation (q/k) & 95.3 & 60.7 & 99.3 & 61.3 & 96.7 & 94.7 & 93.3 & 54.7 & 84.0 & 57.3 \\
NPI \textit{hech} & 93.3 & 2.7 & 96.0 & 5.3 & 79.3 & 86.0 & 21.3 & 4.0 & 61.3 & 13.3 \\
Suspended affixation & 92.7 & 14.7 & 98.0 & 40.7 & 84.7 & 94.0 & 89.3 & 20.7 & 98.0 & 59.3 \\
Nominalization & 100.0 & 90.7 & 99.3 & 92.0 & 99.3 & 100.0 & 98.0 & 52.7 & 74.7 & 19.3 \\
Person agreement & 99.3 & 51.3 & 100.0 & 52.0 & 87.3 & 96.0 & 83.3 & 51.3 & 99.3 & 46.7 \\
Passives & 94.0 & 50.7 & 98.0 & 60.7 & 86.7 & 94.7 & 81.3 & 18.0 & 98.0 & 54.0 \\
Relative clauses & 98.0 & 70.0 & 98.7 & 66.7 & 98.7 & 98.7 & 97.3 & 56.7 & 95.3 & 44.7 \\
Quantifiers (\textit{har bir}) & 96.0 & 80.7 & 98.7 & 82.0 & 92.0 & 97.3 & 84.0 & 97.3 & 98.0 & 82.7 \\
Numeral + singular & 93.3 & 93.3 & 98.0 & 77.3 & 81.3 & 95.3 & 74.0 & 94.0 & 96.0 & 86.7 \\
Postposition case & 100.0 & 48.0 & 100.0 & 64.7 & 99.3 & 100.0 & 100.0 & 57.3 & 100.0 & 72.0 \\
Possessive agreement & 96.7 & 66.7 & 100.0 & 71.3 & 91.3 & 96.7 & 74.7 & 72.7 & 100.0 & 70.7 \\
Dative allomorphy & 88.0 & 1.3 & 92.7 & 14.0 & 70.7 & 87.3 & 48.7 & 67.3 & 90.7 & 37.3 \\
\midrule
All (2,400) & \textbf{95.7} & \textbf{47.6} & \textbf{98.6} & \textbf{50.5} & \textbf{88.5} & \textbf{94.6} & \textbf{76.9} & \textbf{46.2} & \textbf{87.3} & \textbf{47.3} \\
\bottomrule
\end{tabular}
\caption{Accuracy (\%) by model, alphabet of the text and phenomenon; chance is 50. GF-L / GF-C: Goldfish uzn\_latn\_100mb / uzn\_cyrl\_100mb; BERT: tahrirchi-bert-base (PLL); Qwen: Qwen2.5-0.5B. Goldfish-latn is not scored on reformed Latin (82\% of sentences contain an unknown token).}
\label{tab:accuracy}
\end{table*}


\subsection{Overall}
Table~\ref{tab:accuracy} shows that the Uzbek-only models know these phenomena well in the script they were trained on: tahrirchi-bert reaches 98.6\% and Goldfish-latn 95.7\% on Latin text. Both are near ceiling on most phenomena; the hardest are dative allomorphy (88.0\% for Goldfish-latn, 92.7\% for tahrirchi-bert), binding and NPI licensing (90.7\% and 93.3\% for Goldfish-latn). Of tahrirchi-bert's 34 errors on Latin text, 11 are dative allomorphy, where the ungrammatical form (\ungr{qishloqga}) is a non-word that differs from the correct one by a single consonant. Qwen2.5-0.5B is at chance (50.5\%) with extreme per-phenomenon results (7.3\% on binding, 92.0\% on nominalization); \S\ref{sec:length} shows that these follow sentence length rather than grammar.

\begin{table}[t]
\centering\small
\begin{tabular}{lrrr}
\toprule
Comparison & $\Delta$ & 95\% CI & $p$ \\
\midrule
Matched: GF-L on Latin vs GF-C on Cyrillic & +1.1 & [+0.1, +2.2] & 0.044 \\
GF-L: Latin vs Cyrillic text & +7.2 & [+5.9, +8.6] & 3.6e-23 \\
BERT: Latin vs Cyrillic text & +21.7 & [+20.1, +23.4] & 5.8e-108 \\
BERT: current vs reformed Latin & +11.3 & [+9.9, +12.6] & 1.8e-55 \\
Qwen: Latin vs Cyrillic text & +4.4 & [+2.0, +6.8] & 0.00031 \\
\bottomrule
\end{tabular}
\caption{Paired comparisons on the same 2,400 pairs: accuracy difference (points), bootstrap 95\% CI and McNemar's $p$.}
\label{tab:scripts}
\end{table}


\subsection{Same knowledge, different script}
The two Goldfish models score nearly the same when each reads its own script: 95.7\% vs.\ 94.6\% (Table~\ref{tab:scripts}; +1.1 points, 95\% CI +0.1 to +2.2, $p=0.044$), a small Latin advantage at the edge of significance. No single phenomenon differs significantly after Bonferroni correction; the largest gaps are binding (+9.3) and NPI licensing (+7.3), both $p=0.022$. Grammar knowledge of this kind is thus learnable from either script to almost the same degree.

\subsection{Different script, lost knowledge}
When a model reads a script it was not trained on, accuracy drops sharply. Goldfish-latn loses 7.2 points on Cyrillic text, and tahrirchi-bert loses 21.7 points (98.6\% \textrightarrow{} 76.9\%) although its tokenizer covers every Cyrillic character. The loss is uneven across phenomena (Table~\ref{tab:accuracy}): tahrirchi-bert falls to 21.3\% on NPI licensing and 39.3\% on binding, well below chance, and to 48.7\% on dative allomorphy, but stays at 100\% on postposition case and above 97\% on nominalisation and relative clauses. What these models know about Uzbek grammar is tied to the script, not only to the language.

\begin{table}[t]
\centering\small
\begin{tabular}{lrrr}
\toprule
Model & Latin & Cyrillic & Reformed \\
\midrule
Goldfish-latn & 16.3 & 18.6 & -- \\
Goldfish-cyrl & 34.0 & 13.4 & -- \\
tahrirchi-bert & 16.0 & 30.7 & 19.2 \\
Qwen2.5-0.5B & 28.1 & 36.5 & 26.6 \\
\bottomrule
\end{tabular}
\caption{Mean number of scored tokens per sentence. Reading the other script roughly doubles the token count for tahrirchi-bert.}
\label{tab:tokens}
\end{table}


\subsection{Why the other script fails: fragmentation and a length fallback}
\label{sec:length}
Two observations explain this pattern. First, the other script fragments the input: tahrirchi-bert needs 16.0 tokens per Latin sentence but 30.7 per Cyrillic sentence (Table~\ref{tab:tokens}), so the suffixes that carry the grammatical contrast are split into pieces the model never learned to combine. Second, when grammar is no longer available, the model falls back on a preference for shorter sentences. Table~\ref{tab:length} splits the pairs by whether the ungrammatical member is longer, equal or shorter. On Latin text tahrirchi-bert shows no length effect (99.0\%, 98.0\%, 98.8\%); on Cyrillic it is right on 83.0\% of the pairs whose ungrammatical member is longer but only 68.4\% of those where it is shorter, and on reformed Latin 95.6\% vs.\ 79.7\%. The phenomena that fall below chance are exactly those whose ungrammatical member is shorter: binding (\ungr{meni} instead of \textit{oʻzimni}) and NPI licensing (the negation suffix removed). The same heuristic dominates the small multilingual model throughout: Qwen2.5-0.5B is right on 81.4\% of the pairs whose ungrammatical member is longer and on 37.3\% of those where it is shorter, which accounts for both its chance-level total and its extreme per-phenomenon scores.

\begin{table}[t]
\centering\small
\setlength{\tabcolsep}{4pt}
\begin{tabular}{@{}llrrrr@{}}
\toprule
 & & \multicolumn{3}{c}{Ungrammatical member has} & \\
\cmidrule(lr){3-5}
Model & Text & fewer & same & more & $\Delta$ \\
\midrule
GF-L full & Latin & 89.4 & 97.5 & 98.0 & +8.6 \\
GF-L full & Cyrillic & 78.8 & 92.3 & 95.2 & +16.4 \\
GF-C full & Cyrillic & 90.2 & 97.5 & 99.1 & +8.9 \\
GF-C full & Latin & 13.1 & 57.5 & 85.5 & +72.4 \\
tahrirchi-bert & Latin & 99.1 & 98.5 & 98.4 & -0.7 \\
tahrirchi-bert & Cyrillic & 50.2 & 83.2 & 89.4 & +39.2 \\
tahrirchi-bert & Reformed & 58.8 & 93.3 & 95.9 & +37.0 \\
\bottomrule
\end{tabular}
\caption{Accuracy (\%) by whether the ungrammatical member has fewer, the same number of, or more tokens than the grammatical one under the model's own tokenizer. $\Delta$: accuracy on pairs with more tokens minus accuracy on pairs with fewer; 0 means no length effect.}
\label{tab:length}
\end{table}


\begin{table}[t]
\centering\small
\setlength{\tabcolsep}{3.2pt}
\begin{tabular}{@{}lrrrrrr@{}}
\toprule
 & & & \multicolumn{4}{c}{Change in points} \\
\cmidrule(l){4-7}
Model & Cur. & Ref. & both & gr. & other & none \\
\midrule
TahrirchiBERT & 98.7 & 87.3 & $-$11.6 & $-$55.1 & $-$3.4 & 0.0 \\
Qwen2.5-0.5B & 50.4 & 46.8 & $-$6.4 & $-$27.6 & +2.7 & 0.0 \\
\bottomrule
\end{tabular}
\caption{Current against reformed Latin for the models whose tokenizer covers the new letters: accuracy (\%) on all pairs, and the change in points for pairs whose new letters (\"o, \u{g}, \c{s}, \c{c}) fall in the edited word of \emph{both} members ($n=803$), in the edited word of the \emph{gr}ammatical member only ($n=254$), only in \emph{other} words ($n=1220$), or nowhere ($n=123$).}
\label{tab:reform}
\end{table}


\subsection{The 2026 reform}
On reformed Latin, tahrirchi-bert loses 11.3 points (98.6\% \textrightarrow{} 87.3\%). Table~\ref{tab:reform} decomposes this loss by where the new letters fall. When only the rest of the sentence contains new letters, the cost is small (\textminus3.1). When both members' changed words contain new letters, the model loses 11.4 points: a genuine degradation of its grammatical judgments under the new spelling. But when \emph{only the grammatical member} contains a new letter---all binding pairs (\textit{özimni} vs.\ \ungr{meni}) and most nominalization pairs (\textit{-iş} vs.\ \textit{-gan})---accuracy falls from 99.6\% to 44.2\%, below chance: the model prefers the familiarly spelled ungrammatical sentence. This is an effect of orthographic familiarity, not grammar. By phenomenon (Table~\ref{tab:reformphen} in the appendix), binding falls to 29.3\% and nominalisation to 74.7\%, while possessive agreement, postposition case and passives are unaffected, even in the pairs where both members' edited words contain new letters.

\subsection{Source of the sentences}
Five phenomena draw on both our corpus and uz-books, a book collection that some Uzbek models may have been trained on. Within these phenomena, accuracy hardly differs by source for the Uzbek models (Goldfish-latn 96.6\% on corpus pairs vs.\ 97.1\% on book pairs; tahrirchi-bert 99.2\% vs.\ 98.6\%), so we find no sign that pairs from the book collection are easier for them. \todo{seen/unseen analysis against the training data of the evaluated models, where available}

\subsection{Reproducibility}
When we rescored all models after extending the benchmark from 1,200 to 2,400 pairs, all 1,200 shared pairs received identical judgments in all alphabets (log-probabilities within 0.0002). Per-pair scores therefore do not depend on the rest of the benchmark, which also allowed us to update results after replacing individual pairs by rescoring only those pairs.

\section{Discussion}
\label{sec:discussion}

\paragraph{Script dependence.} Our results separate two questions that are easy to conflate: whether a grammatical contrast is learnable from text in a given script, and whether a model trained in one script can use what it learned in another. For matched models the answer to the first is yes, to almost the same degree for Latin and Cyrillic. The answer to the second is largely no, even for a model whose tokenizer covers the other script; the knowledge is stored in terms of the subword units of the training script and does not survive their fragmentation. For Uzbek NLP, this means that models and tools trained before the 2026 reform will make more grammatical errors on text in the new spelling until they are adapted, and that Latin-trained models should not be assumed to handle Cyrillic text.

\paragraph{What else a minimal pair measures.} Two of our findings concern minimal-pair evaluation in general. First, a pair compares two sentences, and anything that makes one of them more familiar to the model---here, spelling---can decide the comparison: when only the grammatical member contains unfamiliar characters, accuracy measures orthographic familiarity, not grammar. Benchmarks that are transliterated, normalised or affected by spelling variation should report accuracy separately for pairs where the critical word changes on both sides. Second, length differences between the members interact with model competence: a model that knows the grammar is indifferent to them, but a model that does not, or that reads an unfamiliar script, falls back on preferring the shorter sentence. Aggregate accuracy then depends on how many pairs happen to have a shorter ungrammatical member, which can push results below chance. This complements the finding of \citet{basar2025turblimp} that length differences explain little of model performance on Turkish: length matters little for models that know the grammar, and a great deal for those that do not. Reporting accuracy by length category, as in Table~\ref{tab:length}, makes this visible; length-normalised accuracy (Appendix~\ref{app:results}) only partly corrects it.

\paragraph{Ceiling effects.} The Uzbek-only models are close to ceiling on most phenomena in their own script, so UzBLiMP mainly discriminates between models that have and have not learned Uzbek, and between scripts. Harder items, for example longer dependencies or phenomena with more lexical variety at the edited position, would be needed to rank strong Uzbek models.

\todo{Grammar-prompting experiment (decision 19) and its discussion, if included.}

\section{Conclusion}

UzBLiMP provides 2,400 native-reviewed minimal pairs for Uzbek in three alphabets, all derived from attested corpus sentences. Evaluating Uzbek and multilingual models shows that grammar knowledge is learnable from either script but does not transfer between them: reading another script fragments the input and leaves models with a preference for shorter sentences. The 2026 Latin reform costs a Latin-trained model 11 points, much of it through orthographic familiarity rather than grammar. We release the benchmark in all three alphabets, the generators and all model outputs. \todo{repository and dataset links}

\section*{Limitations}

\begin{itemize}
\item \textbf{Converted alphabets.} The Cyrillic and reformed versions are converted from Latin, not written natively. Russian loanwords cannot all be recovered by rule, and only the four main rules of the 2026 reform are applied.
\item \textbf{Review.} Judgments were made by the first author and checked with three other native speakers rather than collected from independent annotators, so we do not yet report inter-annotator agreement. \todo{expert review and an independent agreement sample}
\item \textbf{Models.} We evaluate four small models; the Goldfish models are matched in data size but not in content. \todo{larger models}
\item \textbf{Contamination.} The benchmark is built from public text; 962 pairs come from a book collection that one of the evaluated models may have been trained on. \todo{seen/unseen analysis}
\item \textbf{Coverage.} Sentences come from books, Wikipedia and legal texts and reflect written standard Uzbek; colloquial forms were rejected. Several phenomena rely on closed-class items with little lexical variety.
\end{itemize}

\section*{Ethics statement}

The corpora are public text; pairs keep a source identifier so that sentences can be traced and removed if needed. \todo{licences of the book sources before release} Sentences with personal names after the first word or personal data were excluded. An AI assistant (Claude) was used to write code for the generation, conversion and scoring pipeline, to draft parts of this paper, and as an additional consistency check on the native-speaker judgments for some phenomena. All judgments and the final text are the authors' responsibility.

% \section*{Acknowledgments}  % add after review (camera-ready)

\bibliography{custom}

\appendix

\section{Phenomena}
\label{app:phenomena}

Each phenomenon has 150 pairs. We describe the contrast and the edit that produces the ungrammatical member; Table~\ref{tab:examples} gives one example per phenomenon in current and reformed Latin.

\begin{description}
\item[Anaphor agreement.] The reflexive \textit{oʻz-} carries a possessive suffix agreeing with the subject (\textit{siz \dots\ oʻzingizni}); the edit changes the person of the suffix.
\item[Lexical case.] Some verbs govern a fixed case: \textit{qoʻrqmoq} `fear' the ablative, \textit{ishonmoq} `trust' the dative; the edit replaces it with the accusative.
\item[Ditransitive dative.] The recipient of verbs such as \textit{bermoq} `give' and \textit{koʻrsatmoq} `show' is dative; the edit makes a dative pronoun accusative, giving two accusatives in the clause.
\item[Binding.] An object coreferent with the subject must be the reflexive; the edit replaces \textit{oʻzimni} by the plain pronoun \textit{meni}.
\item[Stem alternation.] Stem-final \textit{q}/\textit{k} become \textit{gʻ}/\textit{g} before vowel-initial possessive suffixes; the edit keeps \textit{q}/\textit{k} (a non-word).
\item[NPI \textit{hech}.] \textit{hech kim} `nobody', \textit{hech qachon} `never' etc.\ require a negated verb; the edit removes the negation suffix \textit{-ma-}.
\item[Suspended affixation.] In coordinated first- or second-person past verbs, the first verb keeps its tense and person ending; the edit replaces it with the bare participle \textit{-gan}.
\item[Nominalisation.] Modal predicates (\textit{kerak}, \textit{mumkin}, \textit{lozim}) take the \textit{-ish} verbal noun; the edit replaces it with \textit{-gan}.
\item[Person agreement.] The verb agrees in person with the subject pronoun; the edit changes the person ending.
\item[Passives.] An agent phrase with \textit{tomonidan} requires a passive verb; the edit removes the passive suffix.
\item[Relative clauses.] A relative clause before its head noun uses the participle \textit{-gan}; the edit replaces it with the finite past \textit{-di}.
\item[Quantifiers.] \textit{har bir} `every' takes a singular noun; the edit adds the plural suffix.
\item[Numeral + singular.] A noun after a numeral is singular; the edit adds the plural suffix.
\item[Postposition case.] \textit{keyin} `after' and \textit{beri} `since' take the ablative; the edit replaces it with the dative.
\item[Possessive agreement.] A possessed noun agrees in person with a genitive pronoun (\textit{mening kitobim}); the edit changes the person of the possessive suffix.
\item[Dative allomorphy.] After stem-final \textit{q}/\textit{k} the dative is \textit{-qa}/\textit{-ka}; the edit uses \textit{-ga} (a non-word).
\end{description}

\begin{table*}[t]
\centering\small
\begin{tabular}{lp{6.6cm}p{6.6cm}}
\toprule
Phenomenon & Current Latin & Reformed Latin (2026) \\
\midrule
Anaphor agreement & U oʻzini asta panaga oladi. / *U oʻzimni asta panaga oladi. & U özini asta panaga oladi. / *U özimni asta panaga oladi. \\
Lexical case & Lekin yolgʻizlikdan qoʻrqaman ham. / *Lekin yolgʻizlikni qoʻrqaman ham. & Lekin yolğizlikdan qörqaman ham. / *Lekin yolğizlikni qörqaman ham. \\
Ditransitive dative & Chunki unga vaziyatni tushuntirgandim. / *Chunki uni vaziyatni tushuntirgandim. & Çunki unga vaziyatni tuşuntirgandim. / *Çunki uni vaziyatni tuşuntirgandim. \\
Binding & Qanaqangi yoʻqotishlardan Siz oʻzingizni saqladingiz? / *Qanaqangi yoʻqotishlardan Siz sizni saqladingiz? & Qanaqangi yöqotişlardan Siz özingizni saqladingiz? / *Qanaqangi yöqotişlardan Siz sizni saqladingiz? \\
Stem alternation (q/k) & Uning yuragi hamon gʻash. / *Uning yuraki hamon gʻash. & Uning yuragi hamon ğaş. / *Uning yuraki hamon ğaş. \\
NPI \textit{hech} & Hech kimga yomonlik qilmaganlar. / *Hech kimga yomonlik qilganlar. & Heç kimga yomonlik qilmaganlar. / *Heç kimga yomonlik qilganlar. \\
Suspended affixation & Men tunda oʻtirdim va ishladim. / *Men tunda oʻtirgan va ishladim. & Men tunda ötirdim va işladim. / *Men tunda ötirgan va işladim. \\
Nominalization & Ertagayoq komendaturaga borish kerak. / *Ertagayoq komendaturaga borgan kerak. & Ertagayoq komendaturaga boriş kerak. / *Ertagayoq komendaturaga borgan kerak. \\
Person agreement & Men uning otilishini talab etaman! / *Men uning otilishini talab etamiz! & Men uning otilişini talab etaman! / *Men uning otilişini talab etamiz! \\
Passives & Avrelian fitnachilar tomonidan oʻldirilgan. / *Avrelian fitnachilar tomonidan oʻldirgan. & Avrelian fitnaçilar tomonidan öldirilgan. / *Avrelian fitnaçilar tomonidan öldirgan. \\
Relative clauses & Birdan kutilmagan hol roʻy berdi. / *Birdan kutilmadi hol roʻy berdi. & Birdan kutilmagan hol röy berdi. / *Birdan kutilmadi hol röy berdi. \\
Quantifiers (\textit{har bir}) & Har bir sportchi bittadan sakrashni bajaradi. / *Har bir sportchilar bittadan sakrashni bajaradi. & Har bir sportçi bittadan sakraşni bajaradi. / *Har bir sportçilar bittadan sakraşni bajaradi. \\
Numeral + singular & Ustiga bitta poʻstak yopilgan. / *Ustiga bitta poʻstaklar yopilgan. & Ustiga bitta pöstak yopilgan. / *Ustiga bitta pöstaklar yopilgan. \\
Postposition case & Ish boʻlgandan keyin shu-da. / *Ish boʻlganga keyin shu-da. & Iş bölgandan keyin şu-da. / *Iş bölganga keyin şu-da. \\
Possessive agreement & Jannat mening poyimda boʻlarmish! / *Jannat mening poyingda boʻlarmish! & Jannat mening poyimda bölarmiş! / *Jannat mening poyingda bölarmiş! \\
Dative allomorphy & Biror yoqqa koʻchib ketmaysizlarmi? / *Biror yoqga koʻchib ketmaysizlarmi? & Biror yoqqa köçib ketmaysizlarmi? / *Biror yoqga köçib ketmaysizlarmi? \\
\bottomrule
\end{tabular}
\caption{One pair per phenomenon in current and reformed Latin (the shortest pair whose reformed spelling differs). The Cyrillic versions are in the released data.}
\label{tab:examples}
\end{table*}


\section{Generation and filtering}
\label{app:generation}

\paragraph{Eligible sentences.} A sentence is eligible if it is a single sentence of the target length that starts with a capital letter and ends with sentence-final punctuation, and contains no digits, quotation marks, brackets, dialogue dashes, URLs, foreign letters (w, c outside ch), Cyrillic, capitalised words after the first (except polite \textit{Siz}), or capital letters inside words (an OCR symptom).

\paragraph{Plausibility.} A new word form is accepted only if it occurs at least twice in the 988,950-form lexicon built from the corpus, except for the two morphophonology phenomena. For stem alternation, only stems whose \textit{q}/\textit{k} form is at least 20 times more frequent than the voiced form are used, which excludes non-alternating stems. For possessive agreement, at most 25\% of the pairs use any one genitive pronoun.

\paragraph{Guards.} Generators skip sentences with more than one candidate site, sentences beginning with a proper noun (TurkicNLP tag), function words that only look like the target (\textit{ham}, bare numerals, pronoun case forms for possessives), dative pronouns that are complements of postpositions (\textit{unga qaramay}), accusative pronouns that could be the object of a nearby participle or verbal noun, \textit{-ib} converbs for passives, and words that TurkicNLP does not tag as nouns where a noun is required. Passive and active forms are matched against the allomorphs of the passive suffix (\textit{-il}, \textit{-in}, \textit{-l}, \textit{-n}), and \textit{-gan} assimilates to \textit{-qan}/\textit{-kan} after \textit{q}/\textit{k}. Sentences already used for the same phenomenon in another source are skipped.

\paragraph{Pairs from uz-books.} The rule-generated pairs built from uz-books were filtered with the same criteria for length, names, quotation marks and dialogue, and kept only if the two members differ in one place.

\paragraph{Selection.} For each phenomenon the accepted pairs are shuffled with a fixed seed and the first 150 form the benchmark; the rest form the reserve.

\section{Alphabet conversion}
\label{app:conversion}

\paragraph{Latin to Cyrillic.} TurkicNLP's transliterator with these corrections: \textit{t}+\textit{s} is \cyr{тс} except in loans where \textit{ts} stands for \cyr{ц} (\textit{politsiya} \textrightarrow{} \cyr{полиция}); a tutuq between \textit{s} and \textit{h} is dropped (\textit{Isʼhoq}); \textit{ye} after a consonant becomes \cyr{ъе} after the prefixes \textit{ob}, \textit{sub}, \textit{ad}, \textit{in}, \textit{kon}, \textit{dez}, \textit{pod}, \textit{raz} (\textit{obyekt} \textrightarrow{} \cyr{объект}) and \cyr{ье} elsewhere (\textit{premyera} \textrightarrow{} \cyr{премьера}); \textit{n}+\textit{gʻ} is not the digraph \textit{ng}; Roman numerals are kept; month names and the loan stems \textit{film} and \textit{kompyuter} receive the soft sign. Checks: no Latin letters left, every pair still differs in one word, and Cyrillic-to-Latin round trips reproduce the original except for loans whose Cyrillic spelling is more informative (\cyr{сентябрь}, \cyr{объект}).

\paragraph{Reformed Latin.} oʻ\textrightarrow ö, gʻ\textrightarrow ğ, sh\textrightarrow ş, ch\textrightarrow ç, applied to all words, including capitalised forms. Other details of the reform (e.g.\ the tutuq and proper names) are not covered.

\section{Additional results}
\label{app:results}

Table~\ref{tab:accnorm} gives length-normalised accuracy, Table~\ref{tab:reformphen} the reform effect by phenomenon, and Table~\ref{tab:reasons} the rejection reasons of the review. Scoring ran on CPU with PyTorch, Transformers and minicons; the full evaluation of four models in three alphabets takes about five hours. \todo{exact software versions and hardware}

\begin{table}[t]
\centering\small
\setlength{\tabcolsep}{3.5pt}
\begin{tabular}{@{}lrrrrrr@{}}
\toprule
 & \multicolumn{2}{c}{Latin} & \multicolumn{2}{c}{Cyrillic} & \multicolumn{2}{c}{Reformed} \\
\cmidrule(lr){2-3}\cmidrule(lr){4-5}\cmidrule(l){6-7}
Model & acc & norm & acc & norm & acc & norm \\
\midrule
GF-L 5MB & 63.2 & 68.4 & 53.5 & 56.1 & -- & -- \\
GF-L 10MB & 73.1 & 75.5 & 55.8 & 61.5 & -- & -- \\
GF-L 100MB & 95.8 & 94.5 & 88.3 & 85.3 & -- & -- \\
GF-L full & 96.7 & 95.2 & 91.7 & 88.1 & -- & -- \\
GF-C 5MB & 45.1 & 51.2 & 61.1 & 66.9 & -- & -- \\
GF-C 10MB & 40.3 & 51.0 & 65.0 & 71.5 & -- & -- \\
GF-C 100MB & 47.6 & 53.6 & 94.7 & 93.5 & -- & -- \\
GF-C full & 49.2 & 55.6 & 97.3 & 95.6 & -- & -- \\
TahrirchiBERT & 98.7 & 98.5 & 76.9 & 76.3 & 87.3 & 89.7 \\
BERTbek & 87.6 & 88.1 & -- & -- & -- & -- \\
UzBERT & -- & -- & 84.0 & 85.0 & -- & -- \\
UzRoBERTa & 93.3 & 93.1 & 96.4 & 96.0 & -- & -- \\
Qwen2.5-0.5B & 50.5 & 57.0 & 46.3 & 49.3 & 47.3 & 53.5 \\
\bottomrule
\end{tabular}
\caption{Raw accuracy (acc) and byte-normalised accuracy (norm; log-probability per UTF-8 byte, the main metric of TatBLiMP), in \%.}
\label{tab:accnorm}
\end{table}

\begin{table}[t]
\centering\small
\setlength{\tabcolsep}{4pt}
\begin{tabular}{lrrrrr}
\toprule
Phenomenon & Both & Gr.\ only & Other & Current & Reformed \\
\midrule
Anaphor agreement & 150 & 0 & 0 & 99.3 & 86.7 \\
Lexical case & 60 & 0 & 90 & 100.0 & 87.3 \\
Ditransitive dative & 0 & 0 & 134 & 99.3 & 98.7 \\
Binding & 0 & 150 & 0 & 100.0 & 29.3 \\
Stem alternation (q/k) & 63 & 13 & 64 & 99.3 & 84.0 \\
NPI \textit{hech} & 64 & 0 & 86 & 96.0 & 61.3 \\
Suspended affixation & 64 & 0 & 78 & 98.0 & 98.0 \\
Nominalization & 62 & 88 & 0 & 99.3 & 74.7 \\
Person agreement & 50 & 0 & 77 & 100.0 & 99.3 \\
Passives & 45 & 0 & 92 & 98.0 & 98.0 \\
Relative clauses & 62 & 0 & 84 & 98.7 & 95.3 \\
Quantifiers (\textit{har bir}) & 19 & 0 & 124 & 98.7 & 98.0 \\
Numeral + singular & 37 & 0 & 107 & 98.0 & 96.0 \\
Postposition case & 50 & 0 & 87 & 100.0 & 100.0 \\
Possessive agreement & 39 & 0 & 96 & 100.0 & 100.0 \\
Dative allomorphy & 54 & 0 & 89 & 92.7 & 90.7 \\
\bottomrule
\end{tabular}
\caption{tahrirchi-bert on current vs reformed Latin by phenomenon, with the number of pairs whose new letters fall in both members' changed words, only in the grammatical member's, or only in other words.}
\label{tab:reformphen}
\end{table}

\begin{table}[t]
\centering\small
\setlength{\tabcolsep}{3.4pt}
\begin{tabular}{@{}lrrrr@{}}
\toprule
Phenomenon & Source & Bad OK & Dialect & Meaning \\
\midrule
Anaphor agreement & 9 & 0 & 2 & 0 \\
Lexical case & 26 & 0 & 8 & 0 \\
Ditransitive dative & 17 & 2 & 6 & 0 \\
Binding & 18 & 1 & 6 & 0 \\
Stem alternation & 24 & 0 & 0 & 0 \\
NPI \textit{hech} & 13 & 0 & 6 & 0 \\
Suspended affixation & 13 & 0 & 2 & 0 \\
Nominalisation & 39 & 4 & 5 & 16 \\
Subject agreement & 12 & 3 & 3 & 0 \\
Passives & 30 & 5 & 3 & 0 \\
Relative clauses & 27 & 33 & 4 & 0 \\
Quantifiers & 37 & 7 & 0 & 1 \\
Numeral + singular & 10 & 3 & 4 & 0 \\
Postposition case & 16 & 0 & 6 & 0 \\
Possessive agreement & 19 & 0 & 1 & 0 \\
Dative allomorphy & 19 & 1 & 0 & 0 \\
\midrule
All & 329 & 59 & 56 & 17 \\
\bottomrule
\end{tabular}
\caption{Reasons for rejecting candidate pairs in the native-speaker review. \emph{Source}: the corpus sentence is unnatural or wrong (mostly OCR errors and typos); \emph{Bad OK}: the ungrammatical member has a grammatical reading; \emph{Dialect}: dialectal or colloquial form; \emph{Meaning}: the edit changes the meaning rather than the grammaticality. 1 further pair was rejected because of a personal name.}
\label{tab:reasons}
\end{table}


\end{document}
